\begin{proof}[Proof of \cref{obj:lem:two-sample-l1-lower}]
Write
\[
\mathcal R_{n,k}:=
\inf_{\widehat L}\sup_{R,S\in\Delta_k}
E_{R,S}\!\left[(\widehat L-\|R-S\|_1)^2\right].
\] % lean: CausalSmith.Stat.DiscreteBudgetvalueCurve.two_sample_l1_lower

We first transfer the published known-\(S\) bound to this two-sample experiment. Given a two-sample estimator \(\widehat L\) and a fixed \(S\), average over its second sample:
\[
\overline L_S(X^n):=
E_{Y^n\sim S^{\otimes n}}\!\left[\widehat L(X^n,Y^n)\right].
\]
Conditional Jensen inequality gives, for every \(R\),
\[
E_R\!\left[(\overline L_S-\|R-S\|_1)^2\right]
\leq E_{R,S}\!\left[(\widehat L-\|R-S\|_1)^2\right].
\]
Taking the relevant suprema and infima therefore yields
\[
\sup_{S\in\Delta_k}\inf_{\widetilde L_S}\sup_{R\in\Delta_k}
E_R\!\left[(\widetilde L_S-\|R-S\|_1)^2\right]
\leq \mathcal R_{n,k}.
\] % lean: CausalSmith.Stat.DiscreteBudgetvalueCurve.twoSampleL1_known_minimax_le

The published bound \citep[Theorem 3 and the unknown-both-distributions reduction following Theorem 6]{JiaoHanWeissman2018}, applied with \(a_0=1\) and \(C_0=2\), consequently supplies constants \(\gamma_{\mathrm{pub}}>0\) and \(k_0\) such that
\[
\mathcal R_{n,k}\geq \gamma_{\mathrm{pub}}\,\frac{k}{n\log n}
\quad\text{if}\quad
k\geq k_0,\qquad n\geq\frac{k}{\log k},
\qquad \log n\leq2\log k.
\]
For \(n\geq2\), the last condition and \(\log(ek)=1+\log k\) imply
\[
\mathcal R_{n,k}\geq
\frac{\gamma_{\mathrm{pub}}}{2}\,\frac{k}{n\log(ek)}.
\] % lean: CausalSmith.Stat.DiscreteBudgetvalueCurve.twoSampleL1_lower_dense_logAlphabet

In particular, enlarging a threshold \(k_D\) if necessary, there is a constant \(\gamma_D>0\) such that, whenever \(k\geq k_D\), \(k\geq2\), \(2\leq n\leq k^2\), and \(n\geq k/\log k\),
\[
\mathcal R_{n,k}\geq
\gamma_D\min\left\{1,\frac{k}{n\log(ek)}\right\}.
\]
Indeed, \(n\leq k^2\) gives \(\log n\leq2\log k\), and the minimum is at most its second argument. % lean: CausalSmith.Stat.DiscreteBudgetvalueCurve.twoSampleL1_lower_large_dense

For the smaller sample sizes, use the same dense bound at the integer
\[
n_\circ(k):=\left\lceil\frac{k}{\log k}\right\rceil.
\]
When \(k\geq3\), the inequalities \(1\leq\log k\leq k-1\) show that \(k/\log k>1\), while
\[
\frac{k}{\log k}\leq n_\circ(k)
<\frac{k}{\log k}+1
\leq\frac{2k}{\log k},
\qquad
n_\circ(k)\leq k^2,
\qquad
\log n_\circ(k)\leq2\log k.
\]
Also \(\log(ek)\leq2\log k\), so \(n_\circ(k)\log(ek)\leq4k\). Applying the dense bound at \(n_\circ(k)\) thus gives, for some \(\gamma_S>0\) and threshold \(k_S\),
\[
\mathcal R_{n_\circ(k),k}\geq\gamma_S
\quad (k\geq\max\{k_S,3\}).
\]
If \(1\leq n<k/\log k\), then \(n\leq n_\circ(k)\). An estimator with \(n_\circ(k)\) observations can use their first \(n\) coordinates, whose law is exactly the \(n\)-sample law. Hence
\[
\mathcal R_{n,k}\geq\mathcal R_{n_\circ(k),k}
\geq\gamma_S.
\] % lean: CausalSmith.Stat.DiscreteBudgetvalueCurve.twoSampleL1_lower_large_sparse

Set \(k_\star:=\max\{k_D,k_S,3\}\). To cover \(2\leq k<k_\star\), consider the two parameter pairs
\[
(R_0,S)=((1,0,\ldots,0),(1,0,\ldots,0)),
\qquad
(R_1,S)=((1/2,1/2,0,\ldots,0),(1,0,\ldots,0)).
\]
Their targets are \(0\) and \(1\). Let \(\omega_0\) be the data point at which every observation in both samples has the first category, and put \(\rho_n:=2^{-n}\). The probabilities of \(\{\omega_0\}\) under the two pairs are \(1\) and \(\rho_n\), respectively. Thus, for every estimator \(\widehat L\), its worst-case risk satisfies
\[
\sup_{R,S\in\Delta_k}E_{R,S}\!\left[(\widehat L-\|R-S\|_1)^2\right]
\geq
\max\!\left\{\widehat L(\omega_0)^2,\,
\rho_n\bigl(\widehat L(\omega_0)-1\bigr)^2\right\}
\geq \frac{\rho_n}{4},
\]
since \(\rho_n\leq1\) and \(x^2+(x-1)^2\geq1/2\) for every real \(x\). Taking the infimum over estimators gives
\[
\mathcal R_{n,k}\geq\frac{\rho_n}{4}.
\] % lean: CausalSmith.Stat.DiscreteBudgetvalueCurve.twoSampleL1_finiteAlphabetLower

Minimax risk decreases as the sample size increases by the same first-coordinate argument. Since \(n\leq k^2\leq k_\star^2\), the preceding bound at \(k_\star^2\) gives
\[
\mathcal R_{n,k}\geq
\mathcal R_{k_\star^2,k}\geq
\frac{2^{-k_\star^2}}4
\qquad(2\leq k<k_\star).
\] % lean: CausalSmith.Stat.DiscreteBudgetvalueCurve.twoSampleL1_minimax_antitone

Finally, choose
\[
c:=\min\left\{\gamma_D,\gamma_S,\frac{2^{-k_\star^2}}4\right\}>0.
\]
For \(k\geq k_\star\), split according as \(n<k/\log k\) or \(n\geq k/\log k\); the sparse and dense bounds, respectively, give the claim. For \(2\leq k<k_\star\), use the finite-alphabet bound. In each case the claimed minimum is at most \(1\). % lean: CausalSmith.Stat.DiscreteBudgetvalueCurve.two_sample_l1_lower
\end{proof}
